% ======================================================================
% 4. přednáška (24. 10. 2022)
% Kolektivní stavy (dokončení), jaderné síly, jaderné potenciály,
% radioaktivita (α, β, γ, vnitřní konverze), jaderné reakce, štěpení, fúze
% ======================================================================
\section{Jaderné síly, radioaktivita a jaderné reakce}
\providecommand{\jcCtyrKruh}[1]{\text{\tikz[baseline=(X.base)]\node[draw,circle,inner sep=0.6pt,font=\scriptsize](X){#1};}}

\strana{P04-1}
\subsection{Kolektivní stavy -- pokračování: rotační stavy}

Na minulé přednášce jsme začali mluvit o \kw{kolektivních stavech} jádra, tedy o excitacích,
na kterých se podílí mnoho nukleonů najednou (na rozdíl od jednočásticových excitací ve
slupkovém modelu, kdy se „přestěhuje“ jediný nukleon do vyšší hladiny). Vedle vibrací povrchu
jádra existuje druhý typ kolektivního pohybu -- \kw{rotace}.

Proč by jádro mělo vůbec rotovat „viditelně“? Kvantově mechanicky nemá smysl mluvit o rotaci
kulově symetrického objektu: otočení koule ji nijak nezmění, takže neexistuje žádný stupeň volnosti,
který by se mohl excitovat. Rotační stavy proto pozorujeme jen u \kw{deformovaných jader}
(typicky protáhlých, tvaru „doutníku“ -- prolate), a to při rotaci kolem osy \emph{kolmé} k ose symetrie,
jak ukazuje náčrtek.

\begin{center}
\begin{tikzpicture}[scale=1]
  % deformované jádro (osa symetrie svisle)
  \draw[thick] (0,0) ellipse (0.55 and 1.1);
  \draw[densely dotted,gray] (0,-1.4) -- (0,1.4) node[above,black,font=\small]{osa symetrie};
  % osa rotace vodorovně
  \draw[dashed] (-1.8,0) -- (2.6,0) node[right,font=\small]{osa rotace};
  % šipka rotace
  \draw[-{Stealth[length=2.2mm]},thick] (1.05,0.55) arc[start angle=90,end angle=-200,x radius=0.18,y radius=0.55];
\end{tikzpicture}
\end{center}

Klasicky má rotující těleso energii \(E = L^2/(2I)\), kde \(L\) je moment hybnosti a \(I\) moment setrvačnosti.
V kvantové mechanice nahradíme \(L^2 \to \hbar^2 J(J+1)\), kde \(J\) je kvantové číslo celkového momentu hybnosti
(v poznámkách značené \(\vec{J}\), resp.\ \(j\)). Odtud plyne, že
\begin{itemize}
  \item \textbf{energetické hladiny se škálují s \(\vec{J}\) jako \(j(j+1)\)}:
  \[
    \vzorec{E_\text{rot} = \frac{\hbar^2}{2I}\, J(J+1)} ,
  \]
  \item \textbf{hladiny} (u sudo-sudých jader): \(0^+,\ 2^+,\ \dots\) (dále \(4^+, 6^+, \dots\)).
\end{itemize}
Proč jen sudá \(J\) a kladná parita? Základní stav sudo-sudého jádra je \(0^+\) a deformované jádro
má symetrii vůči otočení o \(180^\circ\) kolem osy rotace (obě „špičky doutníku“ jsou stejné);
tato symetrie vyřadí lichá \(J\). Charakteristickým znakem rotačního pásu je, že hladiny \emph{nejsou ekvidistantní}
(na rozdíl od vibračních): poměr energií \(E(4^+)/E(2^+) = \frac{4\cdot5}{2\cdot3} = 3{,}33\), zatímco pro
harmonické vibrace by byl \(2\). Z naměřené energie \(E(2^+)\) navíc přímo určíme moment setrvačnosti \(I\),
a tedy i míru deformace jádra (viz Davídkovy slajdy „Vibrační a rotační pásy“).

\subsection{Jaderné síly -- popis}

Nyní se vrátíme k otázce, \emph{co} vlastně drží nukleony v jádře pohromadě. V moderním pohledu se
každá interakce (kromě gravitace) popisuje jako \kw{výměna (virtuální) částice}. Elektromagnetická
interakce mezi dvěma náboji je výměna fotonu; protože foton je nehmotný, Coulombův potenciál
\(V\propto 1/r\) má nekonečný dosah. Jaderné síly mají naopak \emph{krátký dosah} (zhruba \(1\)--\(2\,\mathrm{fm}\)),
a Yukawa (1935) z toho usoudil, že zprostředkující částice musí být hmotná. Dnes víme, že na vzdálenostech
\(\gtrsim 1\,\mathrm{fm}\) je to \kw{pion}. Nukleony si mohou vyměnit neutrální pion \(\pi^0\) (proton zůstane protonem,
neutron neutronem) nebo nabitý pion \(\pi^\pm\) -- pak se proton „změní“ v neutron a naopak:

\begin{center}
\begin{tikzpicture}[scale=1,
    fermion/.style={thick,postaction={decorate},decoration={markings,mark=at position 0.55 with {\arrow{Stealth}}}},
    pion/.style={thick,dashed}]
  % diagram 1: p n -> p n pres pi0
  \draw[thick,-{Stealth}] (0,1.6) -- (1.0,1.2);
  \draw[thick] (1.0,1.2) -- (1.6,1.0);
  \draw[thick,-{Stealth}] (1.6,1.0) -- (2.6,1.25);
  \draw[thick] (2.6,1.25) -- (3.2,1.4);
  \draw[thick,-{Stealth}] (0,-0.6) -- (1.0,-0.1);
  \draw[thick] (1.0,-0.1) -- (1.6,0.1);
  \draw[thick,-{Stealth}] (1.6,0.1) -- (2.6,-0.2);
  \draw[thick] (2.6,-0.2) -- (3.2,-0.4);
  \draw[pion] (1.6,1.0) -- (1.6,0.1) node[midway,right]{\(\pi^0\)};
  \fill (1.6,1.0) circle (1.6pt); \fill (1.6,0.1) circle (1.6pt);
  \node[left] at (0,1.6) {\(p\)}; \node[right] at (3.2,1.4) {\(p\)};
  \node[left] at (0,-0.6) {\(n\)}; \node[right] at (3.2,-0.4) {\(n\)};
  % diagram 2: p n -> n p pres pi+
  \begin{scope}[xshift=5.5cm]
  \draw[thick,-{Stealth}] (0,1.6) -- (1.0,1.2);
  \draw[thick] (1.0,1.2) -- (1.6,1.0);
  \draw[thick,-{Stealth}] (1.6,1.0) -- (2.6,1.25);
  \draw[thick] (2.6,1.25) -- (3.2,1.4);
  \draw[thick,-{Stealth}] (0,-0.6) -- (1.0,-0.1);
  \draw[thick] (1.0,-0.1) -- (1.6,0.1);
  \draw[thick,-{Stealth}] (1.6,0.1) -- (2.6,-0.2);
  \draw[thick] (2.6,-0.2) -- (3.2,-0.4);
  \draw[pion,-{Stealth}] (1.6,1.0) -- (1.6,0.1) node[midway,right]{\(\pi^+\)};
  \fill (1.6,1.0) circle (1.6pt); \fill (1.6,0.1) circle (1.6pt);
  \node[left] at (0,1.6) {\(p\)}; \node[right] at (3.2,1.4) {\(n\)};
  \node[left] at (0,-0.6) {\(n\)}; \node[right] at (3.2,-0.4) {\(p\)};
  \end{scope}
\end{tikzpicture}
\end{center}

Hmotnost pionu a z ní plynoucí dosah síly jsou
\[
  m_\pi = 140\,\mathrm{MeV}, \qquad \frac{\hbar c}{m_\pi c^2} \sim 1{,}4\,\mathrm{fm} .
\]
(Číselně \(\hbar c = 197{,}3\,\mathrm{MeV\,fm}\), takže \(197{,}3/140 \approx 1{,}41\,\mathrm{fm}\).) Veličina
\(\hbar/(m_\pi c)\) je \kw{Comptonova vlnová délka} pionu. Heuristicky: virtuální pion „vypůjčí“ energii
\(\Delta E \approx m_\pi c^2\) a podle relací neurčitosti může existovat jen dobu \(\Delta t \approx \hbar/\Delta E\);
za tu dobu urazí nejvýše \(c\,\Delta t = \hbar c/(m_\pi c^2)\). Tím je dána typická vzdálenost, na kterou jaderná síla působí.

Výměna hmotné částice vede na \kw{Yukawův potenciál} (v přirozených jednotkách \(\hbar=c=1\)):
\[
  \vzorec{V(r) \propto -\frac{1}{r}\, e^{-m_\pi \cdot r}} ,
  \qquad\text{tj. v obvyklých jednotkách}\quad V(r)\propto -\frac{1}{r}\,e^{-r/r_0},\ \ r_0 = \frac{\hbar c}{m_\pi c^2}.
\]
Oproti Coulombovu potenciálu je zde navíc exponenciální tlumení: pro \(r \ll r_0\) se potenciál chová jako \(1/r\),
pro \(r \gg r_0\) rychle vymizí. Znaménko minus vyjadřuje přitažlivost.

\textbf{Proč?} Odkud se tento tvar bere? Potenciál je Fourierovou transformací \kw{propagátoru} vyměňované částice:
\[
  V(r) = \frac{-1}{(2\pi)^3}\int e^{i\vec{k}\cdot\vec{r}}\;
  \underbrace{\frac{4\pi}{k^2+m^2}}_{\text{Greenova funkce (propagátor)}}\; \dd^3\vec{k} .
\]
Výraz \(\frac{4\pi}{k^2+m^2}\) je \kw{Greenova funkce} (propagátor) hmotné částice s hmotností \(m\) v hybnostní
reprezentaci; \(\vec{k}\) je přenesená hybnost.

\begin{vysvetleni}[výpočet Fourierova integrálu]
Integrál spočteme ve sférických souřadnicích s osou \(z\) podél \(\vec{r}\). Úhlová část dává
\(\int e^{ikr\cos\theta}\,\dd\Omega = 4\pi\,\frac{\sin kr}{kr}\), takže
\[
  \frac{1}{(2\pi)^3}\int e^{i\vec{k}\cdot\vec{r}}\frac{4\pi}{k^2+m^2}\,\dd^3\vec{k}
  = \frac{1}{(2\pi)^3}\cdot 4\pi\cdot\frac{4\pi}{r}\int_0^\infty \frac{k\sin kr}{k^2+m^2}\,\dd k
  = \frac{16\pi^2}{8\pi^3 r}\cdot\frac{\pi}{2}e^{-mr} = \frac{e^{-mr}}{r},
\]
kde jsme použili známý integrál \(\int_0^\infty \frac{k\sin kr}{k^2+m^2}\dd k = \frac{\pi}{2}e^{-mr}\) (reziduová věta).
Pro \(m\to0\) (foton) dostaneme \(\frac{4\pi}{k^2} \leftrightarrow \frac{1}{r}\), tedy Coulombův potenciál. Hmotnost
vyměňované částice se tedy projeví právě jako „pól“ propagátoru v bodě \(k^2=-m^2\), který po transformaci dá
exponenciální útlum s délkou \(1/m\). Stejná myšlenka (propagátor \(\sim 1/(q^2 - m^2)\)) se později objeví u slabé
interakce a bosonů \(W, Z\) -- jejich velká hmotnost vysvětluje extrémně krátký dosah slabé síly.
Viz Davídkovy slajdy „Jaderné síly (1), (2)“.
\end{vysvetleni}

\subsection{Zpět k potenciálům}

Popis jaderných sil výměnou pionů je fyzikálně názorný, ale pro praktické výpočty stavů jádra (např.\ ve slupkovém
modelu) se používá \emph{efektivní střední potenciál}, v němž se nukleon pohybuje. Přehled modelů, od nejjednodušších:
\begin{itemize}
  \item \kw{jáma} (pravoúhlá potenciálová jáma) -- nejjednodušší, s ostrým okrajem;
  \item \kw{harmonický oscilátor} -- analyticky řešitelný, dobře reprodukuje první magická čísla (2, 8, 20);
\end{itemize}
\strana{P04-2}
\begin{itemize}
  \item \kw{Yukawův potenciál}
  \[
    V(r) = -\frac{1}{r}\, e^{-\alpha\, m_\pi \cdot r},
  \]
  \pozn{písmeno v exponentu čteno jako \(\alpha\); jde zřejmě jen o konstantu udávající dosah, srov.\ \(V\propto -\frac1r e^{-m_\pi r}\) výše};
  \item \kw{Woodsův--Saxonův potenciál}
  \[
    \vzorec{V(r) = \frac{-V_0}{1+\exp\!\left(\frac{r-R}{a}\right)}} ,
  \]
  kde \(V_0\approx 50\,\mathrm{MeV}\) je hloubka jámy, \(R \approx 1{,}2\,A^{1/3}\,\mathrm{fm}\) poloměr jádra
  a \(a\approx 0{,}5\,\mathrm{fm}\) „difuznost“ (tloušťka povrchu).
\end{itemize}
Proč právě Woodsův--Saxonův tvar? Jaderné síly mají krátký dosah, a proto střední potenciál, který cítí nukleon,
kopíruje \emph{rozložení hustoty} nukleonů \(\rho(r)\). To je uvnitř jádra prakticky konstantní (saturace jaderných sil)
a na povrchu plynule klesá k nule (Fermiho rozdělení z experimentů s rozptylem elektronů):

\begin{center}
\begin{tikzpicture}[scale=1]
  \draw[-{Stealth}] (0,0) -- (4.2,0) node[below]{\(r\)};
  \draw[-{Stealth}] (0,0) -- (0,1.8) node[left]{\(\rho(r)\)};
  \draw[thick,domain=0:4,samples=80] plot (\x,{1.3/(1+exp((\x-2.0)/0.35))});
  \draw[dashed,gray] (2.0,0) -- (2.0,0.65) ; \node[below] at (2.0,0) {\(R\)};
\end{tikzpicture}
\end{center}

Woodsův--Saxonův potenciál je tedy „zaoblená“ pravoúhlá jáma: uvnitř má konstantní hodnotu \(-V_0\), v okolí
\(r=R\) na škále \(a\) přechází k nule. Spolu se spin-orbitální interakcí reprodukuje všechna magická čísla.

\subsection{Radioaktivita}

Teď přejdeme od statických vlastností jader k jejich \emph{přeměnám}. Nestabilní jádro se rozpadá náhodně; nelze
předpovědět, kdy se rozpadne konkrétní jádro, ale pravděpodobnost rozpadu za jednotku času \(\lambda\)
(\kw{rozpadová konstanta}) je pro dané jádro konstantní a nezávisí na jeho „stáří“. Za čas \(\dd t\) se tedy rozpadne
\(\dd N = -\lambda N\,\dd t\) jader, a integrací
\[
  \vzorec{N(t) = N(0)\cdot e^{-\lambda \cdot t}}
  \quad\Longrightarrow\quad
  \frac{\dd N}{\dd t} = -\lambda\cdot N(0)\, e^{-\lambda t} .
\]
Z tohoto zákona odvodíme tři užitečné charakteristiky:
\begin{itemize}
  \item \kw{střední doba života} \(\tau\): \(\lambda e^{-\lambda t}\dd t\) je pravděpodobnost, že se jádro rozpadne
  v intervalu \((t, t+\dd t)\); střední hodnota času rozpadu je
  \[
    \tau = \int_0^\infty t\cdot \lambda\cdot e^{-\lambda t}\,\dd t = \frac{1}{\lambda}
  \]
  (per partes). Za dobu \(\tau\) poklesne počet jader na \(1/e \approx 37\,\%\).
  \item \kw{poločas rozpadu} \(T_{1/2}\): doba, za kterou se rozpadne polovina jader,
  \[
    N(T_{1/2}) = \tfrac12 N(0) \quad\Longrightarrow\quad \vzorec{T_{1/2} = \ln 2\cdot\tau} \approx 0{,}693\,\tau .
  \]
  \item \kw{aktivita} -- počet rozpadů za jednotku času, což je veličina, kterou skutečně měříme detektorem:
  \[
    \vzorec{A(t) \equiv \frac{-\dd N(t)}{\dd t} = \lambda\cdot N(t)} ,
  \]
  jednotkou je \kw{becquerel} \(1\,\mathrm{Bq} = 1\,\text{rozpad/s}\) (starší jednotka \(1\,\mathrm{Ci}=3{,}7\cdot10^{10}\,\mathrm{Bq}\)).
\end{itemize}

\begin{priklad}[odhad -- radioaktivita lidského těla (Davídkovy slajdy)]
Dospělý člověk obsahuje asi \(200\,\mathrm{g}\) draslíku, z toho \(0{,}012\,\%\) je radioaktivní \({}^{40}\mathrm{K}\)
(\(T_{1/2}\approx 1{,}25\cdot10^9\) let). Počet jader \(N \approx \frac{200\cdot1{,}2\cdot10^{-4}}{40}\cdot 6\cdot10^{23}
\approx 3{,}6\cdot10^{20}\), a tedy
\(A = \frac{\ln2}{T_{1/2}}N \approx \frac{0{,}693\cdot3{,}6\cdot10^{20}}{1{,}25\cdot10^9\cdot3{,}15\cdot10^7\,\mathrm{s}} \approx 6\cdot10^3\,\mathrm{Bq}\).
Vidíme, že i velmi dlouhý poločas dává měřitelnou aktivitu, je-li jader dost.
\end{priklad}

Podle toho, co jádro při přeměně vyzáří, rozlišujeme rozpady \(\alpha\), \(\beta\) a přechody \(\gamma\); každý je
způsoben jinou interakcí (silná + elektromagnetická, slabá, elektromagnetická).

\subsubsection{Rozpad \texorpdfstring{\(\alpha\)}{alfa}}

Při rozpadu \(\alpha\) jádro vyzáří jádro helia \({}^4_2\mathrm{He}\) (\(\alpha\)-částici):
\[
  {}^A_Z X \;\to\; {}^{A-4}_{Z-2} Y + \alpha .
\]
Proč zrovna \(\alpha\)-částice, a ne třeba jednotlivý proton? Jádro \({}^4\mathrm{He}\) je mimořádně pevně vázané
(\(B\approx 28{,}3\,\mathrm{MeV}\), tj.\ \(7\,\mathrm{MeV}\) na nukleon), takže jeho vyzáření uvolní nejvíce energie.
Rozpad je energeticky možný, pokud
\[
  \vzorec{M(A,Z) > M(A-4,Z-2) + M({}^4_2\mathrm{He})} ,
\]
tj.\ \(Q_\alpha = [M(A,Z) - M(A-4,Z-2) - M({}^4\mathrm{He})]c^2 > 0\) (lze použít atomové hmotnosti, počet elektronů
se vyrovná). Z Weizsäckerovy formule plyne, že to nastává pro
\[
  A \gtrsim 150 .
\]
Prakticky se však \(\alpha\)-rozpad pozoruje hlavně u těžkých jader (\(A\gtrsim 200\)), protože pro malé \(Q_\alpha\)
je poločas astronomicky dlouhý -- důvod uvidíme níže u tunelového jevu.

\strana{P04-3}
\paragraph{Rozpadové řady.} Při rozpadu \(\alpha\) se \(A\) změní o 4, při rozpadu \(\beta\) se \(A\) nemění vůbec.
Zbytek \(A \bmod 4\) se tedy při všech přeměnách zachovává a těžká radioaktivní jádra se dělí do
\textbf{4 řad}:
\begin{center}
\begin{tabular}{lll}
  \(A = 4n\)   & thoriová    & \({}^{232}\mathrm{Th}\to\dots\to{}^{208}\mathrm{Pb}\)\\
  \(A = 4n+1\) & neptuniová  & \({}^{237}\mathrm{Np}\to\dots\to{}^{209}\mathrm{Bi}\)\\
  \(A = 4n+2\) & uranová     & \({}^{238}\mathrm{U}\to\dots\to{}^{206}\mathrm{Pb}\)\\
  \(A = 4n+3\) & aktiniová   & \({}^{235}\mathrm{U}\to\dots\to{}^{207}\mathrm{Pb}\)
\end{tabular}
\end{center}
(Počáteční a koncová jádra jsou doplněna, viz Davídkovy slajdy „Rozpadové řady“.) Neptuniová řada se v přírodě
již nevyskytuje, protože její nejdéle žijící člen \({}^{237}\mathrm{Np}\) má poločas \(2\cdot10^6\) let, což je
zanedbatelné vůči stáří Země; ostatní řady začínají jádry s poločasy srovnatelnými se stářím Země či delšími.

\paragraph{Tunelový jev.} Proč mají \(\alpha\)-zářiče tak rozdílné poločasy (od \(10^{-7}\,\mathrm{s}\) po \(10^{10}\) let),
ačkoli se energie \(\alpha\)-částic liší jen málo (\(4\)--\(9\,\mathrm{MeV}\))? Podívejme se na potenciál, který \(\alpha\)-částice
cítí. Uvnitř jádra je hluboká jáma od přitažlivých \emph{jaderných sil}, vně převládá odpudivá \emph{Coulombova}
bariéra \(\propto 2Z_Y e^2/r\), vysoká typicky \(\sim 25\)--\(30\,\mathrm{MeV}\). Energie \(E_\alpha\) je mnohem nižší
než vrchol bariéry, takže klasicky by \(\alpha\)-částice jádro nikdy neopustila -- kvantově ale bariérou může
\emph{protunelovat}. Pro výpočet nahradíme realistický potenciál (vlevo) zjednodušeným modelem (vpravo):

\begin{center}
\begin{tikzpicture}[scale=0.9]
  % realny potencial
  \draw[-{Stealth}] (0,-2.2) -- (0,2.0) node[left]{\(V\)};
  \draw[-{Stealth}] (-0.3,0) -- (5.2,0) node[below]{\(r\)};
  \draw[thick] (0.05,-1.9) -- (0.9,-1.9) .. controls (1.4,-1.9) and (1.45,-1.0) .. (1.5,0)
     .. controls (1.6,1.2) and (2.0,1.3) .. (2.4,1.1) .. controls (3.2,0.6) and (4,0.35) .. (5,0.25);
  \draw[red,thick] (0,0.45) -- (3.6,0.45);
  \node[red,above] at (0.45,0.45) {\(E_\alpha\)};
  \draw[-{Stealth}] (3.4,1.6) node[right]{Coulomb} -- (2.6,1.15);
  \draw[-{Stealth}] (2.4,-1.4) node[right]{jaderné síly} -- (1.45,-1.0);
  % sipka k modelu
  \draw[-{Stealth},thick] (6.9,1.2) -- (5.9,1.2);
  % model
  \begin{scope}[xshift=8.0cm]
    \draw[-{Stealth}] (0,-2.2) -- (0,2.0) node[left]{\(V\)};
    \draw[-{Stealth}] (-0.3,0) -- (5.2,0) node[below]{\(r\)};
    \draw[thick] (0,-1.9) -- (1.5,-1.9) -- (1.5,1.2) -- (3.2,1.2) -- (3.2,0) -- (5,0);
    \draw[dotted] (0,1.2) -- (1.5,1.2);
    \node[left] at (0,1.2) {\(V_1\)};
    \node[left] at (0,-1.9) {\(-V_0\)};
    \draw[red,thick] (0,0.55) -- (1.5,0.55);
    \node[red,above] at (0.75,0.55) {\(E_\alpha\)};
    \node[below] at (1.5,-1.9) {\(r=a\)};
    \node[below right] at (3.2,0) {\(r=b\)};
    \node[draw,circle,inner sep=1pt,font=\small] at (0.75,-1.1) {1};
    \node[draw,circle,inner sep=1pt,font=\small] at (2.35,-0.4) {2};
    \node[draw,circle,inner sep=1pt,font=\small] at (4.2,0.6) {3};
  \end{scope}
\end{tikzpicture}
\end{center}

Model: oblast \jcCtyrKruh1 \(0<r<a\) je jáma hloubky \(-V_0\) (vnitřek jádra, \(a\approx R\)), oblast \jcCtyrKruh2 \(a<r<b\) je bariéra
výšky \(V_1 > E_\alpha\), oblast \jcCtyrKruh3 \(r>b\) je volný prostor. Řešíme stacionární Schrödingerovu rovnici pro \(\alpha\)-částici
s hmotností \(m_\alpha\):
\[
  \left\{\frac{-\hbar^2}{2m_\alpha}\Delta + V(\vec{x})\right\}\psi(\vec{x}) = E\,\psi(\vec{x}) .
\]
Pro centrální potenciál separujeme radiální a úhlovou část,
\(\psi(\vec{x}) = \frac{u(r)}{r}\,Y_{lm}(\theta,\varphi)\), a pro funkci \(u(r)\) dostaneme jednorozměrnou rovnici
\[
  \left(\frac{\dd^2}{\dd r^2} + \frac{2m_\alpha\,(E-V(r))}{\hbar^2} - \frac{l(l+1)}{r^2}\right)u(r) = 0 .
\]
Člen \(l(l+1)/r^2\) je odstředivá bariéra, která by tunelování dále ztížila; dále uvažujeme \(l=0\). V jednotlivých oblastech
jsou řešení
\begin{align*}
  \jcCtyrKruh1\quad & u = A\cdot\sin\!\left(\sqrt{\frac{2m_\alpha(E+V_0)}{\hbar^2}}\;r\right), \\
  \jcCtyrKruh2\quad & u = B\cdot\exp\!\left(\kappa\, r\right) + C\cdot\exp\!\left(-\kappa\, r\right),
     \qquad \kappa = \sqrt{\frac{2m_\alpha(V_1-E)}{\hbar^2}}, \\
  \jcCtyrKruh3\quad & u = D\cdot\exp\!\left(i\, k\, r\right), \qquad k = \sqrt{\frac{2m_\alpha E}{\hbar^2}} .
\end{align*}
(V poznámkách jsou odmocniny v oblastech 2 a 3 jen naznačeny; doplnili jsme je.) V oblasti 1 bereme jen \(\sin\),
protože \(u(0)=0\) (jinak by \(\psi = u/r\) divergovala). V oblasti 2 je energie pod bariérou, řešení jsou reálné
exponenciály. V oblasti 3 bereme jen \emph{vycházející} vlnu \(e^{ikr}\) -- popisujeme částici, která jádro opouští.
Konstanty \(A,B,C,D\) se určí ze spojitosti \(u\) a \(u'\) v bodech \(r=a\), \(r=b\).

\kw{Pravděpodobnost průchodu} (transmisní koeficient) je poměr hustoty pravděpodobnosti vně a uvnitř:
\[
  T_B = \frac{\left|D\exp\!\left(i\sqrt{\tfrac{2m_\alpha E}{\hbar^2}}\;r\right)\right|^2}{\left|A\sin\!\left(\sqrt{\tfrac{2m_\alpha(E+V_0)}{\hbar^2}}\;r\right)\right|^2} .
\]

\begin{vysvetleni}[proč je \(T_B\) tak citlivé na energii -- Geigerovo--Nuttallovo pravidlo]
Pro širokou a vysokou bariéru (\(\kappa(b-a)\gg1\)) dominuje v oblasti 2 klesající exponenciála a výsledek je přibližně
\(T_B \approx e^{-2\kappa(b-a)}\). Pro skutečnou Coulombovu bariéru se \(\kappa\) mění s \(r\) a exponent nahradíme integrálem
(WKB):
\[
  T_B \approx \exp\!\left(-2\int_a^b \sqrt{\tfrac{2m_\alpha (V(r)-E)}{\hbar^2}}\,\dd r\right)
  \approx \exp\!\left(-\,\text{konst}\cdot\frac{Z}{\sqrt{E_\alpha}}\right)
\]
(Gamowův faktor). \(\alpha\)-částice uvnitř jádra naráží na stěnu s frekvencí \(f \approx v/(2R) \sim 10^{21}\,\mathrm{s}^{-1}\)
a při každém nárazu má šanci \(T_B\) projít, takže \(\lambda = f\,T_B\). Logaritmováním dostaneme
\(\log T_{1/2} = a_1\frac{Z}{\sqrt{E_\alpha}} + a_2\) -- empirické \kw{Geigerovo--Nuttallovo pravidlo} (1911), které vysvětlil
až Gamow (1928) právě tunelováním. Protože \(E_\alpha\) je pod odmocninou v exponentu, změna \(E_\alpha\) o pár MeV změní
poločas o mnoho řádů. Viz Davídkovy slajdy „Tunelový jev (1), (2)“.
\end{vysvetleni}

\strana{P04-4}
\subsubsection{Rozpad \texorpdfstring{\(\beta\)}{beta}}

Rozpad \(\beta\) je projevem \emph{slabé interakce}: mění se v něm neutron na proton nebo naopak, přičemž \(A\) se nemění.
Základním procesem je rozpad \textbf{volného} neutronu
\[
  n \to p + \bar\nu_e + e^- \qquad(\text{volný}),
\]
který je možný, protože \(m_n > m_p + m_e\) (\(Q = 0{,}78\,\mathrm{MeV}\)); volný neutron má střední dobu života asi
\(880\,\mathrm{s}\). Proton je naopak lehčí než neutron, takže volný proton se nerozpadá -- v jádře ale může přeměna
\(p\to n\) nastat, pokud to umožní rozdíl vazebných energií. V jádře jsou proto \textbf{3 varianty}:
\begin{center}
\renewcommand{\arraystretch}{1.4}
\begin{tabular}{lll}
  \(\beta^-\): & \({}^A_Z X \to {}^{\;\;A}_{Z+1} Y + e^- + \bar\nu_e\) & \((n) \to (p) + e^- + \bar\nu_e\)\\
  \(\beta^+\): & \({}^A_Z X \to {}^{\;\;A}_{Z-1} Y + e^+ + \nu_e\) & \((p) \to (n) + e^+ + \nu_e\)\\
  EC:          & \({}^A_Z X + e^- \to {}^{\;\;A}_{Z-1} Y + \nu_e\) & \((p) + e^- \to (n) + \nu_e\)
\end{tabular}
\end{center}
Závorky \((n)\), \((p)\) znamenají nukleon vázaný v jádře. \kw{EC} (electron capture, záchyt elektronu, často \emph{K-záchyt})
znamená, že jádro pohltí elektron z vlastního atomového obalu, nejčastěji z K-slupky; následuje emise charakteristického
rentgenového záření, jak se díra v obalu zaplní. EC konkuruje \(\beta^+\) a je jedinou možností, když
\(0 < Q < 2m_ec^2\).

Pravděpodobnost (rozpadová konstanta) rozpadu \(\beta\) silně roste s uvolněnou energií:
\[
  \vzorec{\lambda \propto Q^5}, \qquad \left(Q = M(X) - M(Y)\right)
\]
(\kw{Sargentovo pravidlo}). Mocnina 5 pochází z fázového prostoru tří částic v koncovém stavu: integrujeme přes hybnosti elektronu
a neutrina, které se dělí o energii \(Q\).
\pozor{Pozor na hmotnosti:} s \emph{atomovými} hmotnostmi platí \(Q=[M(X)-M(Y)]c^2\) pro \(\beta^-\) i EC
(u EC ještě minus vazebná energie zachyceného elektronu), ale pro \(\beta^+\) je
\(Q = [M(X)-M(Y)-2m_e]c^2\), protože atom \(Y\) má o elektron méně a navíc vzniká pozitron.

\paragraph{Spektrum elektronů.} Protože energie \(Q\) se dělí mezi tři částice (elektron, neutrino a zanedbatelně odražené jádro),
elektron nemá pevnou energii, ale \emph{spojité spektrum} od nuly do maxima \(E_\text{max}\approx Q\):

\begin{center}
\begin{tikzpicture}[scale=1]
  \draw[-{Stealth}] (0,0) -- (5.2,0) node[below]{\(E(e^-)\)};
  \draw[-{Stealth}] (0,0) -- (0,2.4) node[left]{\(\frac{\dd N}{\dd E}\)};
  \draw[thick] (0,1.3) .. controls (0.5,2.2) and (1.3,2.2) .. (2.2,1.5) .. controls (3.0,0.8) and (3.8,0.15) .. (4.4,0);
  \draw (4.4,0.1) -- (4.4,-0.1) node[below]{\(Q\)};
\end{tikzpicture}
\end{center}

Právě spojité spektrum vedlo Pauliho (1930) k postulování neutrina -- při dvoučásticovém rozpadu by elektron měl ostrou energii.
Nenulová hodnota spektra u \(E\to0\) pro \(\beta^-\) je důsledkem coulombovského přitahování elektronu ke kladnému jádru
(u \(\beta^+\) je naopak spektrum u nízkých energií potlačeno).

\paragraph{Hmotnostní paraboly.} Které izobary (jádra se stejným \(A\)) jsou stabilní vůči rozpadu \(\beta\)? Ze semiempirické
(Weizsäckerovy) formule plyne, že pro pevné \(A\) je hmotnost \(M(A,Z)\) \emph{kvadratickou funkcí} \(Z\) (Coulombův a asymetrický člen).
Jádra tedy leží na parabole a rozpady \(\beta\) je posouvají „dolů“ k minimu: z levé větve (přebytek neutronů) rozpadem \(\beta^-\),
z pravé (přebytek protonů) rozpadem \(\beta^+\)/EC.

\begin{center}
\begin{tikzpicture}[scale=0.85]
  % lichá jádra
  \draw[-{Stealth}] (0,0) -- (5.6,0) node[below]{\(Z\)};
  \draw[-{Stealth}] (0,0) -- (0,4.2) node[left]{\(M(A,Z)\)};
  \node at (3.0,4.2) {lichá jádra, \(A=\text{konst}\)};
  \draw[thick,domain=0.5:5.0,samples=60] plot (\x,{0.38*(\x-2.6)^2+0.6});
  \foreach \z in {0.8,1.5,2.2,2.9,3.6,4.3} \fill (\z,{0.38*(\z-2.6)^2+0.6}) circle (2pt);
  \draw[-{Stealth},thick] (0.95,2.15) -- (1.4,1.35) node[midway,right]{\(\beta^-\)};
  \draw[-{Stealth},thick] (4.15,1.95) -- (3.65,1.25) node[midway,above right]{\(\beta^+\)};
  \draw (2.2,0.1) -- (2.2,-0.1) node[below]{\(Z\)};
  \draw (2.9,0.1) -- (2.9,-0.1) node[below]{\(Z{+}1\)};
  % suda jadra
  \begin{scope}[xshift=7.6cm]
  \draw[-{Stealth}] (0,0) -- (5.6,0) node[below]{\(Z\)};
  \draw[-{Stealth}] (0,0) -- (0,4.2) node[left]{\(M\)};
  \node at (3.0,4.2) {\(A=\text{sudé, konst}\)};
  \draw[thick,domain=0.4:5.0,samples=60] plot (\x,{0.38*(\x-2.7)^2+1.4});
  \draw[thick,domain=0.4:5.0,samples=60] plot (\x,{0.38*(\x-2.7)^2+0.4});
  \node[right] at (5.0,{0.38*(2.3)^2+1.4}) {L--L};
  \node[right] at (5.0,{0.38*(2.3)^2+0.4}) {S--S};
  \foreach \z in {1.0,2.6,4.2} \fill (\z,{0.38*(\z-2.7)^2+1.4}) circle (2pt);
  \foreach \z in {0.5,1.8,3.4,4.9} \fill (\z,{0.38*(\z-2.7)^2+0.4}) circle (2pt);
  \draw[-{Stealth},thick] (2.6,1.35) -- (1.9,0.75) node[midway,left]{\(\beta^+\)};
  \draw[-{Stealth},thick] (2.6,1.35) -- (3.3,0.65) node[midway,right]{\(\beta^-\)};
  \end{scope}
\end{tikzpicture}
\end{center}

\begin{itemize}
  \item \textbf{Liché \(A\)} (sudo-lichá i licho-sudá jádra): párovací člen je nulový, všechna jádra leží na \emph{jedné parabole}
  a stabilní je zpravidla právě jeden izobar (ten nejblíže minimu).
  \item \textbf{Sudé \(A\)}: párovací člen \(\pm\delta\) rozdělí jádra na dvě paraboly -- \kw{licho-lichá} (L--L, výše) a \kw{sudo-sudá}
  (S--S, níže). Při rozpadu \(\beta\) se střídají L--L a S--S, takže se „skáče“ mezi parabolami. Licho-liché jádro u minima se může
  proto rozpadat \emph{oběma} způsoby, \(\beta^+\) i \(\beta^-\) (typický příklad \({}^{64}\mathrm{Cu}\)), a mohou existovat i dvě stabilní
  sudo-sudá jádra se stejným \(A\) (mezi nimi je možný jen velmi vzácný dvojitý rozpad \(\beta\)). Stabilních licho-lichých jader je
  jen několik lehkých (\({}^2\mathrm{H}, {}^6\mathrm{Li}, {}^{10}\mathrm{B}, {}^{14}\mathrm{N}\)).
\end{itemize}

\strana{P04-5}
\subsubsection{Přechody \texorpdfstring{\(\gamma\)}{gama}}

Po rozpadu \(\alpha\) nebo \(\beta\) často vzniká dceřiné jádro v \emph{excitovaném} stavu. To se pak deexcituje vyzářením fotonu
(\(\gamma\)-přechod, elektromagnetická interakce); \(Z\) ani \(A\) se nemění, energie fotonu je rozdíl energií hladin
(\(E_\gamma\sim 0{,}1\)--\(\text{několik}\,\mathrm{MeV}\)). Typickým příkladem -- a nejpoužívanějším kalibračním a průmyslovým zdrojem -- je
\({}^{60}\mathrm{Co}\):

\begin{center}
\begin{tikzpicture}[scale=1]
  % Co
  \node at (0.9,3.55) {\({}^{60}\mathrm{Co}\)};
  \draw[thick] (0,3.1) -- (1.9,3.1); \node[left] at (0,3.1) {\(2^+\)};
  \node[right] at (1.95,3.1) {\(E_\gamma = 0{,}06\,\mathrm{MeV}\)};
  \draw[thick] (0,2.5) -- (2.6,2.5); \node[left] at (0,2.5) {\(5^+\)};
  \draw[decorate,decoration={snake,amplitude=1.2pt,segment length=4pt},-{Stealth}] (0.4,3.1) -- (0.4,2.52);
  \node[left] at (0.3,2.8) {\(\gamma\)};
  % Ni
  \draw[thick] (3.5,2.2) -- (5.2,2.2); \node[left] at (3.5,2.2) {\(4^+\)};
  \draw[thick] (3.5,1.2) -- (5.2,1.2); \node[left] at (3.5,1.2) {\(2^+\)};
  \draw[thick] (3.5,0.2) -- (5.2,0.2); \node[left] at (3.5,0.2) {\(0^+\)};
  \node[below] at (4.35,0.15) {\({}^{60}_{28}\mathrm{Ni}\)};
  \draw[decorate,decoration={snake,amplitude=1.2pt,segment length=4pt},-{Stealth}] (4.5,2.2) -- (4.5,1.22);
  \node[right] at (5.3,1.7) {\(E_\gamma = 1{,}17\,\mathrm{MeV}\)};
  \draw[decorate,decoration={snake,amplitude=1.2pt,segment length=4pt},-{Stealth}] (4.9,1.2) -- (4.9,0.22);
  \node[right] at (5.3,0.7) {\(E_\gamma = 1{,}33\,\mathrm{MeV}\)};
  % beta prechody
  \draw[red,thick,-{Stealth}] (2.4,2.5) -- (3.8,2.22) node[midway,above right,font=\small]{\(\beta^-\)};
  \draw[red,thick,-{Stealth}] (1.6,3.1) -- (3.9,1.22) node[pos=0.4,left,font=\small]{\(\beta^-\)};
\end{tikzpicture}
\end{center}

Základní stav \({}^{60}\mathrm{Co}\) (\(5^+\), \(T_{1/2}=5{,}27\) let) se rozpadá \(\beta^-\) na hladinu \(4^+\) v \({}^{60}\mathrm{Ni}\);
ta se kaskádou dvou \(\gamma\)-přechodů \(4^+\to2^+\to0^+\) o energiích \(1{,}17\) a \(1{,}33\,\mathrm{MeV}\) dostane do základního stavu.
Proč se \(5^+\) nerozpadá přímo do \(2^+\) nebo \(0^+\)? Rozpad by vyžadoval velkou změnu momentu hybnosti (\(\Delta J = 3\) resp.\ 5), což
je silně potlačené („zakázané“ přechody) -- preferuje se \(5^+\to4^+\) s \(\Delta J=1\).

Stejné výběrové pravidlo vysvětluje existenci \kw{izomeru} \({}^{60\mathrm{m}}\mathrm{Co}\): stav \(2^+\) jen \(0{,}06\,\mathrm{MeV}\)
(přesněji \(58{,}6\,\mathrm{keV}\)) nad základním stavem \(5^+\). Foton nese moment hybnosti \(L\ge1\) a musí platit
\(|J_i-J_f|\le L\le J_i+J_f\); přechod \(2^+\to5^+\) tak vyžaduje \(L\ge3\), což je velmi pomalé, a stav žije \(10{,}5\) minuty. Část izomerů
se proto stihne rozpadnout \(\beta^-\) přímo na hladinu \(2^+\) niklu.

\subsubsection{Vnitřní konverze}

Co když přechod fotonem není možný vůbec? Foton má spin 1, a proto jediný foton nemůže zprostředkovat přechod \(0\to0\)
(\(L=0\) je vyloučeno). Příkladem je \({}^{72}\mathrm{Ge}\), jehož první excitovaný stav je \(0^+\) (\(\approx 0{,}69\,\mathrm{MeV}\)):

\begin{center}
\begin{tikzpicture}[scale=1]
  \draw[thick] (0,2.4) -- (2.5,2.4) node[right]{\(2^+\)};
  \draw[thick] (0,1.2) -- (2.5,1.2) node[right]{\(0^+\)};
  \draw[thick] (0,0) -- (2.5,0) node[right]{\(0^+\)};
  \node[below] at (1.25,-0.05) {\({}^{72}\mathrm{Ge}\)};
  \draw[decorate,decoration={snake,amplitude=1.2pt,segment length=4pt},-{Stealth}] (1.25,1.2) -- (1.25,0.03);
  \node[right] at (1.35,0.6) {\(\gamma\)};
  \draw[thick,red] (0.85,0.25) -- (1.65,0.95);
  \draw[thick,red] (0.85,0.95) -- (1.65,0.25);
\end{tikzpicture}
\end{center}

Přechod \(0^+\to0^+\) se uskuteční \kw{vnitřní konverzí}: jádro předá excitační energii přímo (elektromagnetickou interakcí,
bez reálného fotonu) jednomu z elektronů atomového obalu, nejčastěji z K-slupky, jejichž vlnová funkce má nenulovou hodnotu i uvnitř jádra.
Elektron je emitován s ostrou energií \(E_e = \Delta E - B_K\) (\(B_K\) je vazebná energie v K-slupce) -- na rozdíl od spojitého spektra
\(\beta\). Následuje emise rentgenového záření či Augerových elektronů. Vnitřní konverze konkuruje \(\gamma\)-přechodům obecně, nejen
u \(0\to0\); významná je u těžkých jader a nízkých energií.

\strana{P04-6}
\subsection{Jaderné reakce}

Zatím jsme se zabývali samovolnými přeměnami; nyní budeme jádra bombardovat částicemi. Jaderné reakce slouží jak ke zkoumání
struktury jader (rozptylové experimenty), tak k výrobě izotopů a k získávání energie.
\begin{itemize}
  \item \textbf{Mnoho typů}, značí se zkratkou \(X(a,b)Y\), např.
  \[
    (n,p):\ X + n \to Y + p,\qquad (n,\gamma),\qquad (\alpha,p),\qquad \dots
  \]
  První písmeno je dopadající částice, druhé vyletující. Reakce \((n,\gamma)\) (radiační záchyt) se používá např.\ k výrobě
  \({}^{60}\mathrm{Co}\) ozařováním \({}^{59}\mathrm{Co}\) neutrony v reaktoru; reakcí \((\alpha,p)\) provedl Rutherford (1919)
  první umělou přeměnu prvků \({}^{14}\mathrm{N}(\alpha,p){}^{17}\mathrm{O}\).
  \item \textbf{Přímé} \(\times\) \textbf{přes složené jádro}:
  \[
    a + X \to Y + b \qquad\qquad a + X \to C \to Y + b
  \]
  U \kw{přímých reakcí} interaguje dopadající částice jen s několika nukleony na povrchu a reakce proběhne za dobu průletu jádrem
  (\(\sim10^{-22}\,\mathrm{s}\)). U reakcí přes \kw{složené jádro} \(C\) (Bohr, 1936) se energie dopadající částice rozdělí mezi všechny
  nukleony, vznikne excitovaný a relativně dlouho žijící (\(10^{-16}\)--\(10^{-14}\,\mathrm{s}\)) systém, který „zapomene“, jak vznikl, a
  rozpadne se statisticky -- způsob rozpadu nezávisí na způsobu vzniku.
  \item \textbf{Energetická bilance}:
  \[
    \vzorec{Q = m_a + m_X - m_b - m_Y}\quad(\cdot c^2)
    \qquad
    \begin{cases}
      Q<0 & \dots\ \text{endotermické},\\
      Q>0 & \dots\ \text{exotermické}.
    \end{cases}
  \]
  \(Q\) je rozdíl klidových energií na začátku a na konci, tj.\ o kolik vzroste kinetická energie. Exotermické reakce
  mohou probíhat i s pomalými částicemi (typicky neutrony); endotermické mají \kw{prahovou energii}: v laboratorní soustavě
  (terčík \(X\) v klidu, nerelativisticky) \(E_\text{prah} \approx -Q\left(1+\frac{m_a}{m_X}\right)\), protože část kinetické
  energie se musí spotřebovat na pohyb těžiště a nemůže se přeměnit na hmotnost.
\end{itemize}

Účinný průřez reakce v závislosti na energii vykazuje ostrá maxima -- \kw{rezonance}:

\begin{center}
\begin{tikzpicture}[scale=1]
  \draw[-{Stealth}] (0,0) -- (5.4,0) node[below]{\(E\)};
  \draw[-{Stealth}] (0,0) -- (0,2.4) node[left]{\(\frac{\dd\sigma}{\dd\Omega}\)};
  \draw[thick] (0.3,1.9) .. controls (0.9,1.0) and (1.5,0.7) .. (2.0,0.85)
     .. controls (2.2,0.95) and (2.3,1.15) .. (2.45,0.8) .. controls (2.6,0.4) and (2.8,0.3) .. (3.3,0.3)
     .. controls (3.6,0.3) and (3.7,0.55) .. (3.85,0.3) .. controls (4.0,0.2) and (4.5,0.2) .. (5.0,0.2);
\end{tikzpicture}
\end{center}

Rezonance odpovídají energiím, při nichž soustava \(a+X\) „trefí“ hladinu složeného jádra \(C\) -- pravděpodobnost jeho vytvoření
je pak řádově vyšší. Tvar rezonance popisuje Breitova--Wignerova formule, šířka \(\Gamma\) souvisí s dobou života hladiny
přes \(\Gamma\tau\approx\hbar\). Klesající pozadí u nízkých energií odpovídá zákonu \(\sigma\propto 1/v\) typickému pro pomalé neutrony.

\subsubsection{Štěpné reakce}

Pro těžká jádra je vazebná energie na nukleon (\(\approx7{,}6\,\mathrm{MeV}\) u uranu) nižší než pro jádra střední hmotnosti
(\(\approx 8{,}5\,\mathrm{MeV}\)). Rozdělení těžkého jádra na dvě středně těžká proto uvolní energii. Jádro \({}^{235}\mathrm{U}\) se po
záchytu (i pomalého) neutronu rozštěpí:
\[
  n + {}^{235}_{\;\,92}\mathrm{U} \to X + Y + k\cdot n + m\,\gamma \qquad (k,m\in\mathbb{N}),
\]
kde \(X, Y\) jsou štěpné fragmenty (typicky nesymetrické, s \(A\approx 95\) a \(140\)), v průměru \(k\approx2{,}4\) neutronu.
Právě to, že vznikají další neutrony, umožňuje \kw{řetězovou reakci}. Uvolněná energie je
\[
  \vzorec{Q \approx 180\,\mathrm{MeV} + 20\,\mathrm{MeV} \approx 200\,\mathrm{MeV}},
\]
kde zhruba \(180\,\mathrm{MeV}\) se uvolní okamžitě (hlavně kinetická energie fragmentů, okamžité neutrony a \(\gamma\)) a asi
\(20\,\mathrm{MeV}\) později při rozpadech \(\beta\) fragmentů, které mají přebytek neutronů (část z toho odnesou antineutrina).
Hrubý odhad: \(236\times(8{,}5-7{,}6)\,\mathrm{MeV}\approx 210\,\mathrm{MeV}\), v souladu s uvedenou hodnotou. Pro srovnání,
chemické reakce uvolní řádově \(\mathrm{eV}\) na atom -- tedy \(10^8\)krát méně.

\begin{priklad}[spotřeba paliva v JE Temelín (Davídkovy slajdy)]
Tepelný výkon jednoho bloku \(3\,\mathrm{GW}\) za rok odpovídá
\(\frac{3\cdot10^9\,\mathrm{J/s}\cdot3{,}15\cdot10^7\,\mathrm{s}}{200\cdot10^6\cdot1{,}6\cdot10^{-19}\,\mathrm{J}}\approx3\cdot10^{27}\) štěpení,
tj.\ \(\approx1{,}2\,\mathrm{t}\) \({}^{235}\mathrm{U}\), a při obohacení \(4{,}25\,\%\) asi \(27\,\mathrm{t}\) uranu ročně.
\end{priklad}

\subsubsection{Fúzní reakce}

Na druhém konci křivky vazebné energie leží lehká jádra: jejich sloučením (\kw{fúzí}) do pevněji vázaného jádra se energie
také uvolňuje, a to na nukleon dokonce více než při štěpení. Překážkou je Coulombova bariéra mezi kladnými jádry, kterou je třeba
překonat (či protunelovat) -- proto fúze vyžaduje teploty \(10^7\)--\(10^8\,\mathrm{K}\).
\begin{itemize}
  \item \kw{tokamak} \textbf{ITER} (budovaný mezinárodní experimentální reaktor ve Francii, plazma udržováno magnetickým polem):
  \[
    {}^2_1\mathrm{H} + {}^3_1\mathrm{H} \to {}^4_2\mathrm{He} + n \qquad (Q = 17{,}6\,\mathrm{MeV}).
  \]
  Reakce D--T má nejvyšší účinný průřez při nejnižších dosažitelných teplotách. Většinu energie (\(14{,}1\,\mathrm{MeV}\)) odnáší
  neutron (lehčí částice při dvoučásticovém rozpadu dostane větší podíl kinetické energie).
\strana{P04-7}
  \item \(p + {}^7_3\mathrm{Li} \to {}^4_2\mathrm{He} + {}^4_2\mathrm{He}\) \quad \((Q = 17\,\mathrm{MeV})\) -- historicky první reakce
  vyvolaná uměle urychlenými částicemi (Cockcroft a Walton, 1932), přesněji \(Q=17{,}35\,\mathrm{MeV}\)
  \pozn{hmotnostní číslo Li v originále přepsáno, čteno jako 7}.
  \item \kw{p--p cyklus} a \kw{CNO cyklus} ve Slunci -- v obou se v součtu čtyři protony přemění na \({}^4\mathrm{He}\)
  (\(4p\to{}^4\mathrm{He}+2e^++2\nu_e\)) s uvolněnou energií \(26{,}7\,\mathrm{MeV}\). Ve Slunci dominuje p--p cyklus, v hmotnějších (teplejších)
  hvězdách CNO cyklus, v němž jádra C, N, O slouží jen jako katalyzátor.
\end{itemize}
